\section{Introduction}
\label{sec:introduction}

% Section 1: Task, applications, and challenge

Feed-forward 3D Gaussian Splatting (3DGS) directly predicts a 3D
Gaussian representation~\citep{kerbl20233d} from input images, avoiding
the costly per-scene optimization of conventional 3DGS.
This efficiency makes it attractive for large-scale 3D reconstruction and
digital content creation.
The ultimate goal is a single model that generalizes reliably across
arbitrary scene scales, camera configurations, and target viewpoints, yet
such broad generalization remains challenging.

% Section 2: Limitation of existing normalization spaces
A prevailing paradigm is to first predict per-pixel depth using a
geometry foundation model~\citep{wang2025vggt,lin2025depth,team2026hy,jiang2026omnix,yu2026infinidepth},
and then predict per-pixel Gaussian attributes based on the resulting
geometry in its \emph{radial-distance-normalized} coordinate space.
\emph{Radial Distance Normalization (RDNorm)} uses the mean radial
distance of reconstructed 3D points to the reference-camera origin as a
scene-level scale, and globally rescales the scene to canonicalize this
statistic. However, RDNorm \textbf{does not} directly canonicalize the spacing between
neighboring Gaussian centers, which we refer to as \emph{Gaussian
spacing}. Since adequate surface coverage requires the Gaussian scale to match the
local Gaussian spacing, variations in this quantity directly induce shifts
in the corresponding learning targets.

As illustrated in Fig.~\ref{fig:gsnorm}, visually similar regions can exhibit
substantially different Gaussian spacings after RDNorm under different
scene scales and camera configurations. Consequently, similar local
observations may correspond to different Gaussian-scale targets across
scenes, introducing ambiguity in scale supervision.
A predictor biased toward smaller Gaussian scales may
under-cover regions with large Gaussian spacing and produce holes,
whereas one biased toward larger scales may cause excessive overlap in
densely sampled regions and lead to over-smoothing, as shown in
Fig.~\ref{fig:ablation}.
This ambiguity in Gaussian-scale supervision poses a major bottleneck to
robust feed-forward 3DGS generalization across varying scene scales and
camera configurations.


\begin{figure}[t]
    \centering
    \includegraphics[width=\textwidth]{figures/method/pipeline_1.pdf}
    \vspace{-0.2in}
    \caption{
        \textbf{Normalization space.}
        (a) \emph{Radial Distance Normalization} canonicalizes a global
        radial statistic but leaves Gaussian spacing varying across scenes
        and camera configurations, resulting in inconsistent
        supervision targets.
        (b) \emph{Gaussian Spacing Normalization} canonicalizes a representative statistic of Gaussian spacing, providing a more consistent spatial reference for Gaussian prediction.
    }
    \vspace{-0.2in}
    \label{fig:gsnorm}
\end{figure}


% Section 3: Our insight
To address this issue, we introduce \textbf{SplatAnything}, a
feed-forward 3DGS framework equipped with \emph{Gaussian Spacing
Normalization (GSNorm)}. The key idea of GSNorm is to derive the
scene-level normalization factor from a representative statistic of
Gaussian spacing, rather than from a global radial statistic. Under the
per-pixel Gaussian prediction paradigm, where each valid image pixel
corresponds to a Gaussian center, perspective projection relates Gaussian
spacing to scene depth and camera intrinsics. Based on this relation,
GSNorm computes Gaussian spacing over all valid observations and uses its
median as a robust scene-level normalization factor. The reconstructed
geometry is then rescaled accordingly, canonicalizing the representative
Gaussian spacing across scenes. This provides a more consistent spatial reference for Gaussian-scale prediction, alleviating ambiguity in supervision caused by variations in
scene scale and camera configuration and thereby improving generalization
across diverse scenes and viewpoints.

% Section 4: Our evaluation
To rigorously assess SplatAnything and prior methods, we further
introduce a cross-dataset benchmark for in-the-wild evaluation.
Unlike conventional protocols~\citep{wang2026tttlrm,xu2025depthsplat,
xu2025resplat,ye2026yonosplat,ren2026tokengs}, which often focus on
training-related domains and interpolated target views, our benchmark
covers diverse datasets, scene scales, image resolutions, and camera
configurations, while explicitly varying input-view sparsity and
emphasizing challenging extrapolated viewpoints. Extensive experiments show that SplatAnything
achieves state-of-the-art performance, with particularly strong gains under large scene and viewpoint shifts.

In summary, our contributions are:

\vspace{-0.05in}

\begin{itemize}
    \item
    We identify ambiguity in Gaussian-scale supervision across varying
    scene scales and camera configurations induced by RDNorm. To address
    this issue, we propose \emph{Gaussian Spacing Normalization (GSNorm)},
    which canonicalizes a robust scene-level statistic of Gaussian spacing
    to provide a more consistent spatial reference for Gaussian prediction.

    \item
    We establish a systematic cross-dataset evaluation benchmark for
    feed-forward 3DGS, spanning diverse scene categories, image
    resolutions, camera configurations, input-view sparsity levels, and
    challenging extrapolated viewpoints. This benchmark enables a
    controlled evaluation of cross-scene and cross-view generalization.
\end{itemize}